\documentclass[10pt,a4paper]{amsart}
\usepackage{amsmath,amssymb,amsthm,graphicx}
\usepackage{verbatim}
\newenvironment{reference}{\comment}{\endcomment}
%\newenvironment{reference}{}{}
\newenvironment{slogan}{\comment}{\endcomment}
\newenvironment{history}{\comment}{\endcomment}
% For commutative diagrams we use Xy-pic
\usepackage[all]{xy}

% We use 2cell for 2-commutative diagrams.
\xyoption{2cell}
% Theorem environments.
%
\theoremstyle{plain}
\newtheorem{theorem}[subsection]{Theorem}
\newtheorem{proposition}[subsection]{Proposition}
\newtheorem{lemma}[subsection]{Lemma}

\theoremstyle{definition}
\newtheorem{definition}[subsection]{Definition}
\newtheorem{example}[subsection]{Example}
\newtheorem{exercise}[subsection]{Exercise}
\newtheorem{situation}[subsection]{Situation}

\theoremstyle{remark}
\newtheorem{remark}[subsection]{Remark}
\newtheorem{remarks}[subsection]{Remarks}

\numberwithin{equation}{subsection}

% Macros
%
\def\lim{\mathop{\mathrm{lim}}\nolimits}
\def\colim{\mathop{\mathrm{colim}}\nolimits}
\def\Spec{\mathop{\mathrm{Spec}}}
\def\Hom{\mathop{\mathrm{Hom}}\nolimits}
\def\Ext{\mathop{\mathrm{Ext}}\nolimits}
\def\SheafHom{\mathop{\mathcal{H}\!\mathit{om}}\nolimits}
\def\SheafExt{\mathop{\mathcal{E}\!\mathit{xt}}\nolimits}
\def\Sch{\mathit{Sch}}
\def\Mor{\mathop{\mathrm{Mor}}\nolimits}
\def\Ob{\mathop{\mathrm{Ob}}\nolimits}
\def\Sh{\mathop{\mathit{Sh}}\nolimits}
\def\NL{\mathop{N\!L}\nolimits}
\def\CH{\mathop{\mathrm{CH}}\nolimits}
\def\proetale{{pro\text{-}\acute{e}tale}}
\def\etale{{\acute{e}tale}}
\def\QCoh{\mathit{QCoh}}
\def\Ker{\mathop{\mathrm{Ker}}}
\def\Im{\mathop{\mathrm{Im}}}
\def\Coker{\mathop{\mathrm{Coker}}}
\def\Coim{\mathop{\mathrm{Coim}}}

% Boxtimes
%
\DeclareMathSymbol{\boxtimes}{\mathbin}{AMSa}{"02}

%
% Macros for moduli stacks/spaces
%
\def\QCohstack{\mathcal{QC}\!\mathit{oh}}
\def\Cohstack{\mathcal{C}\!\mathit{oh}}
\def\Spacesstack{\mathcal{S}\!\mathit{paces}}
\def\Quotfunctor{\mathrm{Quot}}
\def\Hilbfunctor{\mathrm{Hilb}}
\def\Curvesstack{\mathcal{C}\!\mathit{urves}}
\def\Polarizedstack{\mathcal{P}\!\mathit{olarized}}
\def\Complexesstack{\mathcal{C}\!\mathit{omplexes}}
% \Pic is the operator that assigns to X its picard group, usage \Pic(X)
% \Picardstack_{X/B} denotes the Picard stack of X over B
% \Picardfunctor_{X/B} denotes the Picard functor of X over B
\def\Pic{\mathop{\mathrm{Pic}}\nolimits}
\def\Picardstack{\mathcal{P}\!\mathit{ic}}
\def\Picardfunctor{\mathrm{Pic}}
\def\Deformationcategory{\mathcal{D}\!\mathit{ef}}
\usepackage{cleveref}
\crefname{originallemma}{original lemma}{original lemmas}
\Crefname{originallemma}{Original lemma}{Original lemmas}
\makeatletter
\newlabel{lemma-base-change-separated}{{lemma-base-change-separated}{source}{}{sourcepointer-1}{}}
\newlabel{lemma-base-change-separated@cref}{{[originallemma][0][]lemma-base-change-separated}{[1][0][]source}{}{}{}}
\makeatother
\crefname{originallemma}{original lemma}{original lemmas}
\Crefname{originallemma}{Original lemma}{Original lemmas}
\makeatletter
\newlabel{lemma-check-separated-covering}{{lemma-check-separated-covering}{source}{}{sourcepointer-2}{}}
\newlabel{lemma-check-separated-covering@cref}{{[originallemma][0][]lemma-check-separated-covering}{[1][0][]source}{}{}{}}
\makeatother
\crefname{originallemma}{original lemma}{original lemmas}
\Crefname{originallemma}{Original lemma}{Original lemmas}
\makeatletter
\newlabel{lemma-properties-covering-imply-diagonal}{{lemma-properties-covering-imply-diagonal}{source}{}{sourcepointer-3}{}}
\newlabel{lemma-properties-covering-imply-diagonal@cref}{{[originallemma][0][]lemma-properties-covering-imply-diagonal}{[1][0][]source}{}{}{}}
\makeatother
\crefname{originallemma}{original lemma}{original lemmas}
\Crefname{originallemma}{Original lemma}{Original lemmas}
\makeatletter
\newlabel{lemma-separated-implies-morphism-separated}{{lemma-separated-implies-morphism-separated}{source}{}{sourcepointer-4}{}}
\newlabel{lemma-separated-implies-morphism-separated@cref}{{[originallemma][0][]lemma-separated-implies-morphism-separated}{[1][0][]source}{}{}{}}
\makeatother
\crefname{originallemma}{original lemma}{original lemmas}
\Crefname{originallemma}{Original lemma}{Original lemmas}
\makeatletter
\newlabel{lemma-separated-over-separated}{{lemma-separated-over-separated}{source}{}{sourcepointer-5}{}}
\newlabel{lemma-separated-over-separated@cref}{{[originallemma][0][]lemma-separated-over-separated}{[1][0][]source}{}{}{}}
\makeatother
\crefname{originallemma}{original lemma}{original lemmas}
\Crefname{originallemma}{Original lemma}{Original lemmas}
\makeatletter
\newlabel{lemma-trivial-implications}{{lemma-trivial-implications}{source}{}{sourcepointer-6}{}}
\newlabel{lemma-trivial-implications@cref}{{[originallemma][0][]lemma-trivial-implications}{[1][0][]source}{}{}{}}
\makeatother
\crefname{originaltheorem}{original theorem}{original theorems}
\Crefname{originaltheorem}{Original theorem}{Original theorems}
\makeatletter
\newlabel{theorem-quasi-DM}{{theorem-quasi-DM}{source}{}{sourcepointer-7}{}}
\newlabel{theorem-quasi-DM@cref}{{[originaltheorem][0][]theorem-quasi-DM}{[1][0][]source}{}{}{}}
\makeatother
\title{Quasi-DM-ness tested through locally quasi-finite covers over test objects}
\author{The Stacks Project Authors}
\date{stacks-project:stacks-morphisms.tex:local-snapshot-20261012; GNU FDL 1.2 or later}
\begin{document}
\maketitle
\noindent\textit{Editorial note.} All delivered mathematical text is the author's own, copied from the pinned Stacks Project chapter source sources/stacks-algebraic-stacks-morphisms/pinned-original-source.tex; no statement or proof step was rewritten, reordered or removed.\par\medskip
\noindent\textit{Editorial note.} Preamble provenance: the Stacks Project keeps one single shared preamble for every chapter; this card takes it directly from the staged sources/stacks-algebraic-stacks-morphisms/preamble.tex as three verbatim spans (lines 8-12, 14-18 and 154-217), so no derived file was written and nothing was added or redefined. The wrapper class and the base packages are supplied by the assembler.\par\medskip
\noindent\textit{Editorial note.} Bibliography provenance: neither the counted statement nor its proof contains a citation (this lane rescans the whole span before writing the card), so the card carries no bibliography and no entry was reconstructed.\par\medskip
\noindent\textit{Editorial note.} Selection provenance: this card was opened under the file-based control contract (CONTROL.md), with the mechanical pre-card gate run on its statement and proof interval before the card was written.\par\medskip
\noindent\textit{Allowed standard prerequisites (editorial).}\begin{itemize}
\item The characterization of local quasi-finiteness through covers over every test object. Invoked inside the counted proof (source lines 5104-5112). Supplies the covers used to test the condition.
\item The criterion for checking a property of a morphism along a flat finitely presented cover. Invoked inside the counted proof (source lines 5113-5125). Transfers the property between the morphism and its base change.
\item The relation between quasi-DM-ness and local quasi-finiteness of the comparison. Used at the end of the counted proof (source lines 5126-5131). Closes the equivalence.
\end{itemize}\medskip
\ifdefined\hypertarget\hypertarget{sourcepointer-1}{}\fi
\noindent\textit{Source reference.} Original lemma lemma-base-change-separated: sources/stacks-algebraic-stacks-morphisms/pinned-original-source.tex, source lines 703--703. Navigation pointer only: the counted text refers to the Stacks item lemma-base-change-separated, which is an existing labeled result of this same chapter and is not part of this card. The pointer carries the target's own label as its number and its definition line as its span; its content is not reproduced here.\par\medskip
\ifdefined\hypertarget\hypertarget{sourcepointer-2}{}\fi
\noindent\textit{Source reference.} Original lemma lemma-check-separated-covering: sources/stacks-algebraic-stacks-morphisms/pinned-original-source.tex, source lines 727--727. Navigation pointer only: the counted text refers to the Stacks item lemma-check-separated-covering, which is an existing labeled result of this same chapter and is not part of this card. The pointer carries the target's own label as its number and its definition line as its span; its content is not reproduced here.\par\medskip
\ifdefined\hypertarget\hypertarget{sourcepointer-3}{}\fi
\noindent\textit{Source reference.} Original lemma lemma-properties-covering-imply-diagonal: sources/stacks-algebraic-stacks-morphisms/pinned-original-source.tex, source lines 1117--1117. Navigation pointer only: the counted text refers to the Stacks item lemma-properties-covering-imply-diagonal, which is an existing labeled result of this same chapter and is not part of this card. The pointer carries the target's own label as its number and its definition line as its span; its content is not reproduced here.\par\medskip
\ifdefined\hypertarget\hypertarget{sourcepointer-4}{}\fi
\noindent\textit{Source reference.} Original lemma lemma-separated-implies-morphism-separated: sources/stacks-algebraic-stacks-morphisms/pinned-original-source.tex, source lines 1061--1061. Navigation pointer only: the counted text refers to the Stacks item lemma-separated-implies-morphism-separated, which is an existing labeled result of this same chapter and is not part of this card. The pointer carries the target's own label as its number and its definition line as its span; its content is not reproduced here.\par\medskip
\ifdefined\hypertarget\hypertarget{sourcepointer-5}{}\fi
\noindent\textit{Source reference.} Original lemma lemma-separated-over-separated: sources/stacks-algebraic-stacks-morphisms/pinned-original-source.tex, source lines 888--888. Navigation pointer only: the counted text refers to the Stacks item lemma-separated-over-separated, which is an existing labeled result of this same chapter and is not part of this card. The pointer carries the target's own label as its number and its definition line as its span; its content is not reproduced here.\par\medskip
\ifdefined\hypertarget\hypertarget{sourcepointer-6}{}\fi
\noindent\textit{Source reference.} Original lemma lemma-trivial-implications: sources/stacks-algebraic-stacks-morphisms/pinned-original-source.tex, source lines 682--682. Navigation pointer only: the counted text refers to the Stacks item lemma-trivial-implications, which is an existing labeled result of this same chapter and is not part of this card. The pointer carries the target's own label as its number and its definition line as its span; its content is not reproduced here.\par\medskip
\ifdefined\hypertarget\hypertarget{sourcepointer-7}{}\fi
\noindent\textit{Source reference.} Original theorem theorem-quasi-DM: sources/stacks-algebraic-stacks-morphisms/pinned-original-source.tex, source lines 4286--4286. Navigation pointer only: the counted text refers to the Stacks item theorem-quasi-DM, which is an existing labeled result of this same chapter and is not part of this card. The pointer carries the target's own label as its number and its definition line as its span; its content is not reproduced here.\par\medskip
\begin{lemma}
\label{lemma-characterize-quasi-DM}
Let $f : \mathcal{X} \to \mathcal{Y}$ be a morphism of algebraic stacks.
The following are equivalent
\begin{enumerate}
\item $f$ is quasi-DM,
\item for any morphism $V \to \mathcal{Y}$ with $V$ an algebraic space
there exists a surjective, flat, locally finitely presented, locally
quasi-finite morphism $U \to \mathcal{X} \times_\mathcal{Y} V$ where
$U$ is an algebraic space, and
\item there exist algebraic spaces $U$, $V$ and a morphism
$V \to \mathcal{Y}$ which is surjective, flat, and
locally of finite presentation, and a morphism
$U \to \mathcal{X} \times_\mathcal{Y} V$ which is surjective, flat,
locally of finite presentation, and locally quasi-finite.
\end{enumerate}
\end{lemma}

\begin{proof}
The implication (2) $\Rightarrow$ (3) is immediate.

\medskip\noindent
Assume (1) and let $V \to \mathcal{Y}$ be as in (2). Then
$\mathcal{X} \times_\mathcal{Y} V \to V$ is quasi-DM, see
Lemma \ref{lemma-base-change-separated}.
By
Lemma \ref{lemma-trivial-implications}
the algebraic space $V$ is DM, hence quasi-DM. Thus
$\mathcal{X} \times_\mathcal{Y} V$ is quasi-DM by
Lemma \ref{lemma-separated-over-separated}.
Hence we may apply
Theorem \ref{theorem-quasi-DM}
to get the morphism $U \to \mathcal{X} \times_\mathcal{Y} V$
as in (2).

\medskip\noindent
Assume (3). Let $V \to \mathcal{Y}$ and
$U \to \mathcal{X} \times_\mathcal{Y} V$ be as in (3).
To prove that $f$ is quasi-DM it suffices to show that
$\mathcal{X} \times_\mathcal{Y} V \to V$ is quasi-DM, see
Lemma \ref{lemma-check-separated-covering}.
By
Lemma \ref{lemma-properties-covering-imply-diagonal}
we see that $\mathcal{X} \times_\mathcal{Y} V$ is quasi-DM.
Hence $\mathcal{X} \times_\mathcal{Y} V \to V$ is quasi-DM by
Lemma \ref{lemma-separated-implies-morphism-separated}
and (1) holds. This finishes the proof of the lemma.
\end{proof}

\end{document}
